\par\Needspace{4\baselineskip}
\SourceLeaf{020b}
\begin{center}
\pdfbookmark[2]{2.}{article-002}2.
\end{center}
The laws just stated differ from those discovered by our \textsc{Kepler} only in that they are expressed in a form applicable to every kind of conic section and that account is taken of the action of the moving body on the Sun, on which the factor $\BareRoot{(1+\mu)}$ depends. If we regard these laws as phenomena inferred from innumerable and indisputable observations, geometry will show what action
\SourceLeaf{021}the Sun must exert on the bodies moving around it for these phenomena to occur continually. In this way it is found that the Sun acts on the bodies revolving around it just as if an attractive force, whose strength is inversely proportional to the square of the distance, impelled the bodies towards the centre of the Sun. Conversely, if we start by assuming such an attractive force as a principle, those phenomena follow from it as necessary consequences. Here it will suffice to have stated the laws. There is all the less need to dwell here on their connection with the principle of gravitation, since after the great \textsc{Newton} several authors have treated this subject, among them the distinguished \textsc{Laplace}, who dealt with it in his consummate work, \textit{Mécanique Céleste}, in a manner that left nothing further to be desired.

\begin{center}
\pdfbookmark[2]{3.}{article-003}3.
\end{center}

Investigations into the motions of celestial bodies, insofar as they take place in conic sections, by no means require a complete theory of this class of curves. Indeed, a single general equation will suffice as a foundation for everything that follows. It seems most appropriate to choose the very equation that we arrive at as the characteristic equation when we investigate a curve described according to the law of attraction. Thus, if we specify any position of the body in its orbit by its distances $x, y$ from two straight lines drawn in the plane of the orbit and intersecting at right angles at the centre of the Sun, that is, at one of the foci of the curve, and if we further denote the distance of the body from the Sun (always to be taken as positive) by $r$, we shall have a linear equation among $r, x, y$, namely $r+\alpha x+\beta y=\gamma$. Here $\alpha, \beta, \gamma$ denote constant quantities, with $\gamma$ by its nature always positive. If we change the otherwise arbitrary orientation of the straight lines to which the distances $x, y$ are referred, provided they continue to intersect at right angles, the form of the equation and the value of $\gamma$ clearly remain the same, whereas $\alpha$ and $\beta$ take different values. It is evident that the orientation can be chosen so that $\beta$ becomes $=0$ and $\alpha$ is in any case nonnegative. If we then write $\alpha, \gamma$ as $e, p$, respectively, our equation takes the form $r+ex=p$. The line to which the distances $y$ are then referred is called the \emph{line of apsides}, $p$ the \emph{semiparameter}, and $e$ the \emph{eccentricity}. Finally, the conic section is called an \emph{ellipse}, a \emph{parabola}, or a \emph{hyperbola}, according as $e$ is less than, equal to, or greater than unity.

\SourceLeaf{022}It is also readily seen that the position of the line of apsides is completely determined by the conditions given, with the single exception of the case in which both $\alpha$ and $\beta$ were already $=0$. In that case we always have $r=p$, whatever lines the distances $x, y$ are referred to. Since then $e=0$, the curve (which will be a circle) belongs to the class of ellipses by our definition. Its distinctive feature is that the position of the apsides remains entirely arbitrary, if indeed we extend that notion to this case as well.
